\chapter{Terms}

\section{The words this paper uses}

\textbf{Purpose:} Say what each technical word in this paper means, using the name the field already gives it. \textbf{Scope:} this paper and the Salishan paper, \texttt{PUBLIC/\allowbreak{}delta\_\allowbreak{}null/\allowbreak{}} and \texttt{PUBLIC/\allowbreak{}Salishan/\allowbreak{}} in theory\_bucket; the engine workbook, \texttt{theory/\allowbreak{}workbook/\allowbreak{}} in orior at commit \texttt{1044ca6}; and the tools under \texttt{maint/\allowbreak{}data/\allowbreak{}} in the same tree.

Most of the words here are standard and need no explanation. This chapter covers the rest: words with a narrow meaning in this paper, and ideas whose field name a reader might not expect. Where a field has a name for an idea, the paper uses that name and cites where it comes from. A file path or a function name in the code keeps its own name, because the path has to find the file.

\subsection{Sets and arrangements}

\noindent\textbf{alphabet.} The set of symbols a pattern and a text are written in, written \(\Sigma\). The paper assumes nothing about it: no order, no arithmetic, and no promise that it is finite or can be listed. When the structure of the data defines a topology on it, that topology is written \(\tau_\Sigma\).

\noindent\textbf{configuration.} A finite arrangement of points or offsets. It is the shape of one pattern when order does not matter. The alphabet supplies the symbols a pattern could hold, and the configuration says which positions the pattern uses.

\noindent\textbf{family.} A collection whose members are themselves sets or configurations, such as every choice of probe positions, or every shuffled copy of the data.

\noindent\textbf{orbit.} Everything one configuration becomes under a stated group of moves, such as rotations or translations. Two configurations share an orbit only once the group and how it acts are stated.

\noindent\textbf{residue class.} A congruence class modulo a stated number. It is never used here to mean a bucket, a cluster or a grid cell.

\noindent\textbf{subconfiguration.} A subset \(A \subseteq D\) of a pattern's positions, used as a filter. Checking only these positions keeps the pattern's structure and drops some of its conditions.

\noindent\textbf{support.} The finite set of positions a pattern occupies, written \(D \subset G\), or the set of outcomes a distribution gives nonzero probability. Both are the standard meanings of the word.

\subsection{The permutation null}

\noindent\textbf{permutation test.} A test that builds its null distribution from the data itself, by shuffling it many times and computing the same statistic on every shuffled copy \cite{src:Good-2005}. Here the shuffles keep the count of every symbol and any stated grouping. If \(\Pi_\Sigma\) is the set of those shuffles, then for a statistic \(M\) the quantity this paper reports is
\[
\Delta_\Sigma[M] = M(x) - \mathbb{E}_{\pi \sim \Pi_\Sigma}\!\left[M\bigl(\pi(x)\bigr)\right],
\]
the observed statistic minus its mean over the shuffled copies. It is a difference from a null and not a new operator. It says what the arrangement of the data adds beyond its symbol counts, and it does not say what produced the arrangement.

\noindent\textbf{maximum entropy reference.} The shuffled copies are the most disordered arrangements the symbol counts allow. Building a reference that way is Jaynes's maximum entropy principle \cite{src:Jaynes}. The reference uses the counts and assumes nothing else. Entropy is strictly concave and a count constraint is linear. Exactly one distribution meets the constraint, and the reference is unique.

\noindent\textbf{free energy above equilibrium.} The distance from a measured state to its own maximum entropy state under a constraint. Thermodynamics gives that quantity this name. The engine workbook (orior \texttt{1044ca6}, \texttt{theory/\allowbreak{}workbook/\allowbreak{}chapters/\allowbreak{}chapter\_\allowbreak{}orior\_\allowbreak{}workbook.\allowbreak{}tex:57}, \texttt{:1218}) records the match as exact and cites Parrondo, Horowitz and Sagawa for it \cite{src:Parrondo}.

\noindent\textbf{relative entropy.} Also called Kullback-Leibler divergence and written \texttt{D(p‖q)}, a distance between two distributions \cite{src:Kullback}. The match to free energy holds for the whole reference and does not carry over to each statistic. Section 4.3.1 of the first chapter shows that what decides a probe order's worth is not a relative entropy, and the engine workbook records the same (\texttt{chapter\_\allowbreak{}orior\_\allowbreak{}workbook.\allowbreak{}tex:57}, \texttt{:1218}). The tools compute either a spread of gaps against shuffled copies, or a standardized residual against pooled counts. Both measure a difference from a maximum entropy reference, and neither is a relative entropy.

\noindent\textbf{standardized residual against a pooled null.} The test in \texttt{maint/\allowbreak{}data/\allowbreak{}salishan/\allowbreak{}orior\_\allowbreak{}algorithmic\_\allowbreak{}extraction/\allowbreak{}boundary\_\allowbreak{}check.\allowbreak{}py} that pools the counts of every run, spreads them evenly, and scores each run as its own binomial. A standardized residual is the observed count less the expected, over the square root of the expected \cite{src:Haberman-1973}. The code calls the function \texttt{radix}. It has nothing to do with a number base.

\noindent\textbf{rarer half, commoner half.} The symbols split at the median of their frequencies. Every difference from the null reported in the workbook is over the rarer half unless it says otherwise.

\subsection{Searching}

\noindent\textbf{filter and verify.} The search this whole paper rests on. Check a few positions of the pattern first, drop any place in the text that fails one of them, and run the full comparison only on the places left. A filter that can never drop a true match is called lossless, and filtering followed by verification is standard in string matching \cite{src:Navarro}. Here the positions to check can be any subset of the pattern, chosen in any order.

\noindent\textbf{probe.} One position of the pattern checked by the filter. A place in the text whose symbol at that offset differs from the pattern's cannot be a match. A probe is a necessary condition for a match and never a sufficient one.

\noindent\textbf{probe sequence.} Probes applied one after another, each one checking only the places the previous ones kept. The depth \(\log_2 N / H_2\) appears throughout, and means the number of probes an independence model predicts before one candidate is left. The engine places at most four. How many it places depends on the number asked for, the candidates available and the text (orior \texttt{1044ca6}, \texttt{src/\allowbreak{}engine/\allowbreak{}c/\allowbreak{}engine/\allowbreak{}orior.h:39}, \texttt{:738-743}; \texttt{orior.c:1003-1006}, \texttt{:1019-1023}).

\noindent\textbf{candidate.} A place in the text that every probe so far has kept. It can still turn out not to match.

\noindent\textbf{false negative.} A true match the filter drops. A lossless filter has none, and every bench row in this paper counts them.

\noindent\textbf{expected comparisons before a mismatch.} How far into a pattern a mismatch is expected to appear, \(m(1-2^{-H_2})\), for a pattern of length \(m\).

\noindent\textbf{equality test.} In the engine, the function the caller supplies to say whether two positions hold the same symbol (orior \texttt{1044ca6}, \texttt{src/\allowbreak{}engine/\allowbreak{}c/\allowbreak{}engine/\allowbreak{}orior.h:354-357}). The engine never reads a symbol any other way.

\noindent\textbf{test oracle.} In the corpus work, a hand extraction of what a paper holds, used as the answer a tool is graded against. This is the software testing term \cite{src:Howden-1978}. The hand extractions also work the other way: they define what is known. Anything a tool finds outside them belongs to what is not yet known. Written languages with long recorded rule sets, such as Chinese, put scene and subject into one unit and state the writer's meaning directly. The conjecture is that such languages are easier to sort by family and by author, and that tracing why could show a wider boundary. Nothing here measures that.

\subsection{Corpora}

\noindent\textbf{reference corpus.} The corpus a target is compared against, in the corpus linguistics sense. In classification it is the negative class, and in signal detection it is the background.

\noindent\textbf{gold standard corpus.} A corpus checked by hand so that it can be trusted as correct \cite{src:Wissler-2014}. Here it is the output holding only speech in the target language, with reconstructions, rejected forms, other languages and analysis lines removed. It is checked in two directions against a hand extraction, and \texttt{PUBLIC/\allowbreak{}Salishan/\allowbreak{}chapters/\allowbreak{}chapter\_\allowbreak{}Salishan\_\allowbreak{}pure\_\allowbreak{}corpus\_\allowbreak{}README.\allowbreak{}tex}, generated by orior's \texttt{maint/\allowbreak{}data/\allowbreak{}salishan/\allowbreak{}hand\_\allowbreak{}extraction/\allowbreak{}pure\_\allowbreak{}corpus\_\allowbreak{}index.\allowbreak{}py} (theory\_bucket \texttt{README.md:64-\allowbreak{}68}), lists the files in it.

\noindent\textbf{extraction error classes.} The kinds of damage a PDF text extraction does, recorded for each paper to trace an error to its source. There are five: substitution, collapse, transposition, insertion and deletion. \texttt{maint/\allowbreak{}data/\allowbreak{}salishan/\allowbreak{}corpus\_\allowbreak{}script\_\allowbreak{}extraction/\allowbreak{}repairs.\allowbreak{}py} repairs the ones that can be repaired.

\noindent\textbf{split-half estimate.} The distance between two halves of one corpus, used as the smallest difference a comparison between corpora has to beat. Splitting one set of measurements in two and comparing the halves is the split-half method \cite{src:Spearman-1910, src:Brown-1910}. Cutting a corpus at its midpoint measures the order its papers are in, not the resolution, and \texttt{theory/\allowbreak{}Salishan/\allowbreak{}chapters/\allowbreak{}chapter\_\allowbreak{}Salishan\_\allowbreak{}refs.\allowbreak{}tex} shows it. The halves are therefore sampled by alternating.

\noindent\textbf{byte-pair histogram.} A text reduced to the counts of each pair of adjacent bytes, by the function \texttt{squash} in \texttt{src/\allowbreak{}engine/\allowbreak{}python/\allowbreak{}instrument/\allowbreak{}orior.\allowbreak{}py} at orior \texttt{d09b489}. In the field this is a bigram distribution, the digram approximation of Shannon's paper \cite{src:Shannon}. A broader idea sits behind it, and nothing here measures it. An object spread over many dimensions cannot be known in full without measuring it, and nothing measures the infinite directly. What can be observed is how far the object is from its own most disordered state, because that state is fixed by what the object is made of. Flattening the object to a set of counts keeps that distance and loses the rest, and the counts are never cut off at a bound, because the object's own makeup already sets its limits.

\subsection{Measurements}

\enlargethispage{4\baselineskip}

The two measurements below are different, and a row measured with one cannot stand in for a row measured with the other. Both rest on Proposition 1.

\noindent\textbf{the permutation null.} The difference between an arrangement and the shuffled copies of the same symbols. Each copy is a Fisher-Yates shuffle \cite{src:Fisher} (orior \texttt{src/\allowbreak{}python/\allowbreak{}engine/\allowbreak{}analysis/\allowbreak{}measure/\allowbreak{}period.\allowbreak{}py:23}). It carries most of the results in the workbook, and no outside source has the answer to check it against. The conjecture behind it, which nothing here measures, is that the structure in \(\Sigma\) shuffles what the null set \(\text{null}_0\) could carry, and that the shuffled order shows what happens to \(\text{null}_0\).

\noindent\textbf{the shift agreement detector.} Finds a period or an offset by how often the data agrees with a shifted copy of itself. Published cell edges in the Crystallography Open Database \cite{src:Crystallography-Open-Database} check it three separate ways.

The engine is \texttt{src/\allowbreak{}engine/\allowbreak{}c/\allowbreak{}engine/\allowbreak{}orior.h} and \texttt{orior.c} in orior at commit \texttt{1044ca6} (\texttt{docs/\allowbreak{}steering.md:4}). The bench rows in the first chapter were measured with an earlier library whose modules have Latin names, \texttt{impensa\_\allowbreak{}ancorae\_\allowbreak{}acus} among them (Section 9 of that chapter).

\textbf{Author:} dstroy0 (Douglas Quigg) <dquigg123@gmail.com> \textbf{Date:} 2026-09-04
